% SPDX-License-Identifier: CC-BY-4.0
% Copyright (c) 2026 Martin Schr\"oder <info@swedishembedded.com>

\section{Introduction}
\label{sec:intro}

A transformer language model \citep{vaswani2017attention} treats the $d$
coordinates of its hidden state as unrelated numbers: every map that reads or
writes the residual stream is a dense learned matrix, and permuting the
coordinates, with the matching permutation of the weights, changes nothing.
The feed-forward sublayers hold most of the parameters and arithmetic of a
small model, and at width $d$ they cost $\Theta(d^2)$ per token.

Fourier methods replace dense maps with structured ones elsewhere in deep
learning. Fourier neural operators learn per-frequency weights for operators
on functions \citep{li2020fno}; FNet and adaptive Fourier neural operators mix
tokens along the sequence with Fourier transforms \citep{leethorp2021fnet,
guibas2021afno}; FourierFT stores fine-tuning updates as a few spectral
coefficients of the weight change \citep{gao2024fourierft}. None of these makes
the hidden state of a language model itself a spectrum. Doing so is tempting:
a hidden state of $m$ Fourier coefficients of a longer virtual signal is
cheap to store and cache, and per-frequency operators are cheap to apply.

The obstacle is that a basis change by itself does nothing.
\Cref{prop:absorption} shows that a decoder whose maps on the residual stream
are dense and learned absorbs any orthogonal basis into its weights, and that
truncating the spectrum of such a model yields an ordinary narrow model. A
spectral language model can only differ from a conventional one through
operators that are defined relative to the basis.

\paragraph{Scope.} This paper is about language models whose hidden state is
a truncated real Fourier spectrum and whose feed-forward mixing is done by
per-mode filters around a signal-domain nonlinearity: the representation, the
operators and their gradients, the training and decoding path, the spectral
structure of a conventional model's hidden states, and matched comparisons
against conventional transformers. It does not study Fourier token mixing
along the sequence axis or spectral compression of weights.

\paragraph{Question.} At a fixed hidden width, does a feed-forward mixer that
is diagonal in a Fourier basis around a signal-domain gated product reach the
language-modeling quality of a dense feed-forward network, and at what
parameter and arithmetic cost?

\paragraph{Contributions.}
\begin{itemize}
\item The basis-absorption argument (\cref{prop:absorption}), which identifies
  the controls a spectral model must beat: a conventional decoder at the
  spectral model's retained width, not at its virtual signal length.
\item SpectralLM (\cref{sec:method}): a decoder with an $m$-coefficient
  spectral residual stream and a $C$-channel spectral mixer with $3Cm$
  parameters, implemented in an existing training and inference engine with
  gradient-checked operators, a fast-transform path gated against a dense
  reference, and incremental decoding (\cref{sec:implementation}).
\item A registered comparison (\cref{sec:protocol}) against width-matched
  conventional decoders trained identically, against the pretrained Qwen3-0.6B
  \citep{qwen3report} on the same held-out text in bits per byte, and a
  measurement of how much of Qwen3-0.6B's hidden-state energy and predictive
  information a truncated Fourier basis retains compared with principal
  components and random orthogonal bases.
\end{itemize}
